\SourceLeaf{416}%
all have the same signs as the sines of \(\vartheta,\vartheta',\vartheta''\), respectively.

I must also mention that in singular cases the general solution given for the problem either ceases to apply altogether or requires some modification. Here I confine myself to indicating the most important points.

I. If one of the observed angles is \(=0\) or \(=180^0\), it is advantageous to take the corresponding observation point as \((1)\) or \((3)\), even if it lies between the other two. If we choose the latter, all parts of the general solution remain valid. We need only regard \(b\) as infinitely large and retain it as a symbol in the calculations. It then drops out of all the results on its own. Clearly, however, the angle cannot be \(0\) or \(180^0\) at more than one point. This occurs only when all three lie on the line \((4),(5)\), in which case the problem would be indeterminate.

II. If two of the observed angles are equal, then two of the points \((1^*),(2^*),(3^*)\) coincide, or one of the quantities \(a,b,a+b\) becomes \(=0\), and hence \(fF=0\) as well. The cubic equation then becomes a quadratic equation. It is also easy to see that the first coefficient of the cubic equation vanishes \textit{only} when two angles are equal.

III. The given solution cannot be applied if the quantities \(A,B\) are proportional to \(a,b\), so that \(\lambda=0\). In this case the cubic equation has been brought in improperly and contains a root extraneous to the problem, while the correct root occurs twice. It is clear that the two combinations by which we derived from \([2]\) and \([3]\) the equations \([4]\) and \([5]\) are then \textit{not distinct}. The equations are therefore identical and each on its own gives \(x=\frac A{2a}=\frac B{2b}\). Evidently we must still use one of the two equations \([2],[3]\). Combining it with \([1]\) readily gives \(cx=(t-n)\BareRoot{(1-xx)}\). It follows that line \(\mathrm{(I)}\) passes through either point \((4)\) or point \((5)\), and
